\begin{table*}[t]
\centering
\small
\caption{\textbf{Example of human-guided revision.}
An A2 candidate that fails surface leakage (H3) in human validation and its revision. Place names are fictional. The gold answer is A in both versions.}
\label{tab:example_humanrev}
\begin{tabular}{@{}p{0.07\textwidth}p{0.43\textwidth}p{0.43\textwidth}@{}}
\toprule
 & \textbf{Korean} & \textbf{English translation} \\
\midrule
Before
& 은빛군은 군 명칭과 군청을 그대로 두고 있지만, 행정구역 개편 후 현재 관내 주소가 `청람광역시 은빛군'으로 표시된다. 다만 일부 행정 안내문에는 아직 `백산도 은빛군'이라는 이전 표기가 남아 있다. 은빛군 관내의 통행료를 받지 않는 도로에 설치된 안전표지를 새로 교체하려 한다. 이 안전표지의 도로교통법상 원칙적 설치·관리권자는 누구인가?
\newline A) 청람광역시장이다.
\newline B) 청람광역시 은빛군수이다.
\newline C) 백산도지사이다.
\newline D) 백산도 은빛군수이다.
& Eunbit County keeps its name and county office, but since an administrative reorganization, addresses within it are written ``Eunbit County, Cheongram Metropolitan City.'' Some administrative notices, however, still carry the former notation ``Eunbit County, Baeksan Province.'' A safety sign on a toll-free road in Eunbit County is to be replaced. Under the Road Traffic Act, who is in principle responsible for installing and managing this sign?
\newline A) The mayor of Cheongram Metropolitan City.
\newline B) The head of Eunbit County, Cheongram Metropolitan City.
\newline C) The governor of Baeksan Province.
\newline D) The head of Eunbit County, Baeksan Province. \\
\midrule
Annotator (H3)
& 지문이 현재 주소는 ``청람광역시 은빛군''이고 ``백산도 은빛군''은 ``아직 남아 있는 이전 표기''라고 밝혀, 법령 지식 없이 읽기만으로 `백산도'가 들어간 C와 D가 지워진다.
& The stem states that the current address is ``Eunbit County, Cheongram Metropolitan City'' and that ``Eunbit County, Baeksan Province'' is a former notation that still remains, so C and D, which contain ``Baeksan Province,'' are eliminated by reading alone, without knowledge of the law. \\
\midrule
After
& 은빛군은 군청 명칭과 군수 직함을 계속 사용하지만, 관내 법정 주소는 `청람광역시 은빛군'으로 표시된다. 은빛군 안의 통행료를 받지 않는 도로에 설치된 노후 안전표지를 교체하려 한다. 이 안전표지의 도로교통법상 원칙적 설치·관리권자는 누구인가?
\newline A) 청람광역시장이다.
\newline B) 은빛군수이다.
\newline C) 청람광역시경찰청장이다.
\newline D) 청람광역시 시설관리공단장이다.
& Eunbit County continues to use the name of its county office and the title of county head, but the legal address within it is written ``Eunbit County, Cheongram Metropolitan City.'' An old safety sign on a toll-free road in Eunbit County is to be replaced. Under the Road Traffic Act, who is in principle responsible for installing and managing this sign?
\newline A) The mayor of Cheongram Metropolitan City.
\newline B) The head of Eunbit County.
\newline C) The commissioner of the Cheongram Metropolitan Police Agency.
\newline D) The director of the Cheongram Metropolitan City Facilities Management Corporation. \\
\bottomrule
\end{tabular}
\end{table*}